\documentclass[11pt,a4paper]{article}
\usepackage[utf8]{inputenc}
\usepackage[T1]{fontenc}
\usepackage[margin=2.3cm]{geometry}
\usepackage{amsmath,lmodern,booktabs,array,listings,xcolor,hyperref,enumitem}
\hypersetup{colorlinks=true,linkcolor=blue!50!black,urlcolor=blue!50!black}
\lstset{basicstyle=\ttfamily\footnotesize,breaklines=true,frame=single,
  keywordstyle=\color{blue!60!black},commentstyle=\color{green!40!black},language=C++}
\newcommand{\gpio}[1]{\texttt{GPIO#1}}
\title{Flapping Butterfly Clip\\\large Build and Concepts Guide (ESP32-C3 + 0.42'' OLED)}
\author{Electronics-AndresForero}
\date{\today}

\begin{document}
\maketitle
\tableofcontents

\section{What you are building}
A hair clip with butterfly wings that flap when you touch it. The same firmware runs in two phases:
\begin{description}[leftmargin=2.2cm,style=nextline]
  \item[Phase 1, bench] USB powered, push-button, any 5\,V micro servo. Goal: make the code and mechanism work.
  \item[Phase 2, wearable] Single-cell LiPo, 2\,g servo, TTP223 touch module, battery monitor. Goal: small, light, safe.
\end{description}
The ESP32-C3 board has a built-in 72$\times$40 pixel OLED that shows the state, battery voltage and servo angle.

\paragraph{Credit.} Inspired by the \emph{Flapping Butterfly Clip} by \textbf{@resetwithregine} on TikTok
and the \emph{Flapping Butterfly Clip v2 Upgrades} build guide on Notion, who also published the v1 gear
mechanism. The firmware, CAD and this document are an independent rewrite; no STL files, code or images from
the original are reused. Please support the original creator, and credit them and this repository if you build on this work.

\section{Safety first}
\begin{itemize}
  \item Use a LiPo with a \textbf{protection PCB}. Never use a bare cell.
  \item Add a \textbf{hard power switch}. Stopping the wings does \emph{not} disconnect the battery.
  \item Charge \textbf{off the clip}, at or below the cell's rated current, on a non-flammable surface, never unattended.
  \item Insulate every joint with heat-shrink. Keep electronics off the metal clip and your skin.
  \item Enclose the gears. Keep hair and fingers out of the mechanism; disconnect the battery before changing wings.
  \item Check the boost converter output with a multimeter \emph{before} connecting the board.
\end{itemize}

\section{Bill of materials}
The same list, with more detail, lives in \texttt{hardware/BOM.md}. Items marked \emph{verify} are typical
values, not measurements of your parts: check them with calipers or a multimeter.

\subsection{Phase 1, bench prototype}
Everything is already in the workshop; nothing needs to be bought.
\begin{center}
\begin{tabular}{@{}p{5.2cm}cp{6.3cm}@{}}
\toprule
Item & Qty & Notes \\ \midrule
ESP32-C3 board + 0.42'' OLED & 1 & 72$\times$40 px, USB-C \\
Micro servo 5\,V (MG90S or SG92R) & 1 & MG90S first; SG92R is untested \\
Push-button (or a jumper to GND) & 1 & trigger on \gpio{4} \\
Electrolytic capacitor 470\,\textmu F & 1 & across \texttt{V5}/\texttt{GD} \\
Breadboard, header strip, wire & -- & \\
USB-C cable, PLA/PETG filament & -- & \\
M2 self-tapping screws & 3 & horn + two flange screws (supplied with the servo) \\
\bottomrule
\end{tabular}
\end{center}

\subsection{Phase 2, wearable: items you already own}
\begin{center}
\begin{tabular}{@{}p{5.2cm}cp{6.3cm}@{}}
\toprule
Item & Qty & Notes \\ \midrule
ESP32-C3 board + 0.42'' OLED & 1 & same board as Phase 1 \\
Boost converter MT3608 & 1 & set to 5.0\,V before connecting; optional, see below \\
Resistors 100\,k$\Omega$ & 2 & battery divider \\
Electrolytic capacitor 100--470\,\textmu F, $\ge$10\,V & 1 & smallest can that fits \\
TP4056 charger & 1 & charge \emph{off} the clip only \\
Heat-shrink tubing & few\,cm & insulate every joint \\
Soldering station, multimeter & -- & \\
\bottomrule
\end{tabular}
\end{center}

\subsection{Phase 2, wearable: items to buy}
\begin{center}
\begin{tabular}{@{}p{5.2cm}cp{6.3cm}@{}}
\toprule
Item & Qty & Notes \\ \midrule
Micro servo, \textbf{2\,g, 3.7--4.2\,V} & 1 & the weight-critical part; fallback 3.7\,g servo \\
1S LiPo, 150--300\,mAh, \textbf{with protection PCB} & 1 & e.g. 501230 / 502030 class; check connector polarity and pitch \\
TTP223 capacitive touch module & 1 & 3.3\,V, digital output \\
Slide switch (SS12D00G) & 1 & hard battery cut-off \\
Silicone hook-up wire 28--30\,AWG & $\sim$1\,m & flexible and light \\
Cable ties 2.5\,mm & 2 & strap the case to the clip through the floor slots \\
\bottomrule
\end{tabular}
\end{center}

\subsection{Printed parts}
Six parts in \texttt{stl/}, source in \texttt{cad/butterfly.scad}: drive gear, idler gear, left wing,
right wing, case tray and case lid (PLA or PETG). The CAD uses no magnets; wings slot into the gear posts.
The case is sized for an SG92R/MG90S class servo; the 2\,g servo needs the \texttt{micro2g} preset re-measured.

\subsection{Power notes}
\begin{itemize}
  \item \textbf{With the boost:} battery $\to$ switch $\to$ MT3608 (5\,V) $\to$ board \texttt{V5}. The board stays
        stable down to about 3.0\,V on the cell; costs about 2\,g, one part and roughly 15\,\% loss.
  \item \textbf{Without the boost:} with a 3.7--4.2\,V servo, feed the cell straight to \texttt{V5} and raise
        \texttt{VBAT\_CUTOFF\_MV} to 3500, because the board's regulator drops out near 3.5\,V.
  \item \textbf{Runtime (estimate, not measured):} roughly 150--200\,mA from the cell, so about 1--1.5\,h from
        250\,mAh. Measure it in Phase 1.
  \item \textbf{Firmware cutoff:} wings stop at about 3.3\,V (3.5\,V without the boost); the cell's own
        protection only cuts near 2.5\,V.
\end{itemize}

\section{Concepts}
\subsection{How a hobby servo is controlled}
A servo expects a pulse every 20\,ms (50\,Hz). The \emph{width} of that pulse sets the angle:
about 0.5\,ms is 0$^\circ$ and about 2.4\,ms is 180$^\circ$ (values in \texttt{config.h}). The ESP32Servo
library generates this with the chip's hardware PWM (LEDC), so it costs no CPU time.

\subsection{Easing}
Jumping straight to a new angle is jerky. We move along
\[ e(t) = \tfrac12 - \tfrac12\cos(\pi t), \qquad t\in[0,1], \]
which accelerates then decelerates. The angle is $\theta = \theta_0 + (\theta_1-\theta_0)\,e(t)$.

\subsection{Non-blocking code}
\texttt{delay()} freezes the program. Instead, every loop iteration reads \texttt{millis()}, computes where
the wing \emph{should} be at that instant, and moves on. That is why touch, battery and display keep working during a flap.

\subsection{Debouncing}
Mechanical contacts chatter for a few milliseconds. A new level is accepted only after it stays unchanged for
\texttt{DEBOUNCE\_MS}, and one event is produced per press.

\subsection{Voltage divider and the ADC}
A LiPo reaches 4.2\,V but the ADC tolerates about 2.5\,V. Two equal resistors give
$V_{pin}=V_{bat}\,R_2/(R_1+R_2)=V_{bat}/2$. The firmware multiplies the reading by 2. Calibrate with a multimeter
using \texttt{VBAT\_CAL}.

\subsection{Why a boost converter (and when to skip it)}
The board's own regulator needs about 3.5\,V or more to hold 3.3\,V. As the cell sags, the boost converter keeps
the board stable down to about 3.0\,V and lets you use more of the capacity, at the price of one extra part,
about 2\,g, and roughly 15\,\% conversion loss. With a 3.7--4.2\,V servo you may skip it, connect the cell
straight to \texttt{V5}, and raise \texttt{VBAT\_CUTOFF\_MV} to 3500.

\subsection{Servo supply rule}
A 3.7--4.2\,V servo is powered from the \emph{battery}; a 5\,V servo from the \emph{boosted} rail. Connecting
the wrong one either under-drives the servo or damages it. Grounds are always common.

\subsection{Low-battery latch}
Under load the cell voltage dips. The firmware low-pass filters the reading, requires several consecutive low
samples, and then \emph{latches} the fault until reboot, so the wings do not restart the moment the voltage recovers.
The cell's own protection PCB cuts at about 2.5\,V, far too low to be relied upon.

\section{Wiring}
\textbf{Check the pins first}: flash \texttt{hwcheck} and confirm the OLED answers at \texttt{0x3C} on
\gpio{5}/\gpio{6}.
\begin{center}
\begin{tabular}{@{}lll@{}}
\toprule
Signal & Pin & Note \\ \midrule
Servo signal & \gpio{3} & \\
Trigger & \gpio{4} & button to GND (P1) / TTP223 OUT (P2) \\
Battery sense & \gpio{0} & ADC1, Phase 2 \\
OLED SDA/SCL & \gpio{5}/\gpio{6} & on-board, do not use \\
Avoid & \gpio{2}, \gpio{8}, \gpio{9} & boot strapping pins \\
\bottomrule
\end{tabular}
\end{center}

\subsection{Phase 1}
Servo red to \texttt{V5}, brown to \texttt{GD}, signal to \gpio{3}; button between \gpio{4} and \texttt{GD};
470\,\textmu F across \texttt{V5}/\texttt{GD}.

\subsection{Phase 2}
\begin{enumerate}
  \item Battery $+$ $\to$ slide switch $\to$ node \texttt{VBAT}. Battery $-$ is common GND.
  \item \texttt{VBAT} $\to$ MT3608 \texttt{IN+}. Set the output to 5.0\,V with a multimeter, then \texttt{OUT+} $\to$ board \texttt{V5}.
  \item \texttt{VBAT} $\to$ servo red (3.7--4.2\,V servo). Servo brown to GND. Signal to \gpio{3}.
  \item \texttt{VBAT} $\to$ 100\,k $\to$ \gpio{0} $\to$ 100\,k $\to$ GND.
  \item TTP223: VCC to \texttt{V3}, GND to GND, OUT to \gpio{4}.
  \item 100--470\,\textmu F across the MT3608 output.
\end{enumerate}

\section{Firmware walkthrough}
\begin{tabular}{@{}ll@{}}
\toprule
File & Role \\ \midrule
\texttt{include/config.h} & every pin and tunable, build-variant switches \\
\texttt{include/WingServo.h} & eased, non-blocking flapping; detaches when idle \\
\texttt{include/TouchButton.h} & debounced press detection for both input types \\
\texttt{include/Battery.h} & divider reading, filter, percent, latched low flag \\
\texttt{include/StatusDisplay.h} & four-line OLED status \\
\texttt{src/main.cpp} & state machine \\
\texttt{tools/hwcheck/main.cpp} & I2C scan and ADC check \\
\bottomrule
\end{tabular}

\paragraph{State machine.}
\texttt{IDLE} $\xrightarrow{\text{touch}}$ \texttt{FLAPPING} $\xrightarrow{\text{touch}}$ \texttt{IDLE};
any state $\xrightarrow{\text{low battery}}$ \texttt{LOW\_BATT} (latched, Phase 2).

\paragraph{The core of the flap} (from \texttt{WingServo.h}):
\begin{lstlisting}
float t = float(now - t0_) / float(FLAP_MS);   // 0..1 progress of this stroke
float e = 0.5f - 0.5f * cosf(PI * t);          // sine ease in-out
angle_  = from_ + (to_ - from_) * e;           // interpolate between the ends
servo_.write(int(angle_ + 0.5f));
\end{lstlisting}

\section{Build, flash, test}
\begin{lstlisting}[language=bash]
pio run -e hwcheck -t upload && pio device monitor   # confirm pins
pio run -e bench    -t upload                        # Phase 1
pio run -e wearable -t upload                        # Phase 2
\end{lstlisting}
\begin{enumerate}
  \item \textbf{hwcheck:} an I2C device at \texttt{0x3C} on \gpio{5}/\gpio{6}. If not, fix \texttt{config.h} before continuing.
  \item \textbf{bench, no servo:} OLED shows IDLE; each press toggles FLAPPING and the angle number moves.
  \item \textbf{bench, with servo:} wings flap smoothly; if gears bind, lower \texttt{ANGLE\_CLOSED\_DEG}.
  \item \textbf{wearable, on a bench supply or the 2000\,mAh cell:} compare the displayed voltage with a multimeter, then set \texttt{VBAT\_CAL}.
  \item \textbf{Cutoff:} run from a supply you can lower slowly; confirm LOW BATT appears near the cutoff and wings stop.
  \item \textbf{Runtime:} measure the real current. The estimate (1--1.5\,h from 250\,mAh) is unverified.
\end{enumerate}

\section{Troubleshooting}
\begin{tabular}{@{}p{4.8cm}p{9.5cm}@{}}
\toprule
Symptom & Likely cause \\ \midrule
Board resets when the servo moves & weak supply; add/enlarge the capacitor, shorten wires \\
OLED blank & wrong pins; run \texttt{hwcheck} \\
Servo buzzes at the end stops & travel too large; reduce \texttt{ANGLE\_CLOSED\_DEG} \\
Battery reads wrong & divider values or \texttt{VBAT\_CAL}; check with a multimeter \\
Touch toggles by itself & TTP223 floating or sharing a long wire; keep it short and grounded \\
\bottomrule
\end{tabular}

\end{document}
